\begin{proposition}[Genus-one curves with nonreduced conductor]\hypertarget{tag-02WL}{}\label{night:P4.14}\label{stmt:P4.14}
\StatementRecord{night:P4.14}{proved}{02WL}{the definitions and constructions stated in the proof}

The reducible-curve inference in \cite[proof of Proposition~7.3,
p.~16 of the attached-source edition]{CDP20} omits the two plane cubics
\[
 Z(Y^2-XZ)=0,\qquad XY(X-Y)=0.
\]
They are reduced Gorenstein curves of arithmetic genus one with
trivial dualizing sheaf and rational normalization, but have
respectively two tangent components and three concurrent
components. Their conductor order is two on every normalization
branch. These curves refute the asserted nodal-cycle
classification at that step.
\end{proposition}
\begin{proof}

\begin{enumerate}[leftmargin=*]
\item Plane adjunction gives $\omega_C=\OO_C(3-3)=\OO_C$ and
$p_a(C)=1$. The first cubic has a line tangent to a smooth
conic at $[1:0:0]$. Its local equation is $y(y-x^2)=0$.
The normalization restrictions are $(p(t,0),p(t,t^2))$.
Their image is exactly the pairs $(a,b)$ with $b-a$ divisible
by $t^2$: the inverse expression is
$p(x,y)=a(x)+y(b(x)-a(x))/x^2$. Thus the normalization
quotient has length two and the conductor exponent is two on
each branch.

\item For the second cubic the branches are $(t,0),(0,t),(t,t)$.
All constant coefficients must agree, and the third linear
coefficient is the sum of the first two. These three conditions
are sufficient. In every degree $k\ge2$, the monomials
$x^k-x^{k-1}y$, $y^k-xy^{k-1}$ and $x^{k-1}y$ independently
realize the coefficients on the three branches. Hence the
normalization quotient has length three and each conductor
exponent is again two. The normalization Euler characteristics
are two and three; subtracting these quotient lengths gives
$\chi(\OO_C)=0$ in both cases.

\item Let $C_2$ and $C_3$ be the two cubics above. The normalization
of $C_2$ is the disjoint union of two copies of $\Pp^1$, parametrized
by $[s:t]\mapsto[s:t:0]$ on the line and
$[s:t]\mapsto[s^2:st:t^2]$ on the conic. The parameters $t/s$ agree
at the two preimages of the tacnode. Exchanging these copies with
these parameters preserves their local ring
$\{(a,b):b-a\in(t^2)\}$ and descends to an involution $\sigma_2$
of $C_2$. In the local coordinates above it is
$(x,y)\mapsto(x,x^2-y)$. The projective automorphism
\[
 \sigma_3[X:Y:Z]=[-Y:X-Y:Z]
\]
has order three and cyclically permutes the components of $C_3$.
For $r=2,3$, take $E'=\C/(\Z+i\Z)$ and let $\tau_r$ be
translation by $1/r$. Its order is exactly $r$, and every nonidentity
power acts freely. Thus the diagonal action defines the compact
surface and its map to an elliptic curve
\[
 X_r=(C_r\times E')/\langle(\sigma_r,\tau_r)\rangle
       \longrightarrow E=E'/\langle\tau_r\rangle.
\]
The cyclic group acts transitively on the components of
$C_r\times E'$, so $X_r$ is irreducible. Let $\widetilde\sigma_r$
be the lift to the normalization $\nu_r:C_r^\nu\to C_r$ induced
by the displayed parameter maps, so
$\nu_r\widetilde\sigma_r=\sigma_r\nu_r$. The normalization of $X_r$ is
$(\coprod_{j=1}^r\Pp^1\times E')/\langle(\widetilde\sigma_r,\tau_r)\rangle$.
Choose a component $C_{r,0}$ of $C_r^\nu$. Restriction of the quotient
map to $C_{r,0}\times E'$ is an isomorphism. Its inverse sends a class
represented by $(z,e)$ with $z\in\widetilde\sigma_r^j C_{r,0}$ to
$(\widetilde\sigma_r^{-j}z,\tau_r^{-j}e)$; this expression is unchanged
on each orbit and is holomorphic on every component.
The reduced inverse-image conductor is the single section
$\Gamma'=\{p\}\times E'$, its conductor divisor is $2\Gamma'$, and
its map to the singular curve $E$ is the finite unramified cover
$E'\to E$ of degree $r$. Analytically locally on $E$, the image
conductor has algebra $\OO_E[u]/(u^2)$ for $r=2$ and
$\OO_E[u,v]/(u,v)^2$ for $r=3$. These local algebras follow from
the normalization constraints in the first two steps and remain
nonreduced after the free quotient. The constructions give a single
horizontal conductor component of multiplicity two with generic
degree two or three over the image conductor.

\end{enumerate}
\end{proof}